\subsection{Repository} \label{I.1}
Before even thinking of editing ay of the documents, create a branch of the repository to work on. The branch should be named "Changes - semester" (Ex: Changes - F26).
The branch should exist and main repository should remain unchanged until all edits have been verified (see \ref{I.3}).

\subsection{Changelog} \label{I.2}
During the lifetime of your editing branch(\ref{I.1}), a changelog must exist, ideally a google/word document with each change made.
I highly suggest separating changes to different documents, and noting the section or subsection that the change takes place.
The changelog should be accessible via the chapter google drive.

\subsection{Editing Rules} \label{I.3}
The editing rules for the documents is as follows:

\begin{enumerate}
    \item \subsecitem{a}{Formatting}
    All changes to any of the documents should follow the formatting detailed in Part \ref{II}. 
    \item \subsecitem{b}{Secondary Inspection}
    All changes should be double checked by yourself and ideally another for errors.
    \item \subsecitem{c}{Branch Merging}
    The branch(\ref{I.1}) should only be merged after:
    \begin{enumerate}
        \item Any document that has been edited has been read through for any missed errors or broken latex.
        \item The relevant \textbf{Last Revised} dates have been updated.
        \item Bylaws Specific: All changes in the changelog(\ref{I.2}) have been inputted into the \textbf{Bylaws Amendment Report} to the Association.
    \end{enumerate}
\end{enumerate}